% ============================================================================
% UGAF-ITS PAPER REVISION — LaTeX Content for Insertion
% ============================================================================
%
% This file contains ready-to-paste LaTeX content for:
%   1. NEW Section VII-B: Computational Validation via Prototype Tool
%   2. NEW Table: Multi-Scenario Comparative Results
%   3. REVISED abstract sentences
%   4. REVISED Section VII-D (Threats to Validity) updates
%
% Instructions:
%   - Insert Section VII-B between current VII-A and VII-B (renumber VII-B→VII-C etc.)
%   - Replace current Table 9 with the new multi-scenario table
%   - Update abstract with the new sentences
%   - Update VII-D with the revised text
% ============================================================================


% ============================================================================
% 1. NEW SECTION VII-B: COMPUTATIONAL VALIDATION
% ============================================================================
% Insert AFTER Section VII-A (Case Study Context) and BEFORE the current
% Section VII-B (UGAF-ITS Application). Renumber subsequent subsections.

\subsection{Computational Validation via Prototype Tool}
\label{subsec:prototype_tool}

To validate the operational feasibility of UGAF-ITS beyond the single
Dubai-inspired scenario, we developed a Python-based governance engine
that algorithmically evaluates ITS deployment architectures against the
unified control catalog. The tool ingests a YAML deployment descriptor
defining the system's AI components, their tier placement, risk levels,
and ownership, and produces a compliance report through five computations:
(i) tier-aware control activation, (ii) evidence backbone construction
with cross-framework reuse analysis, (iii) bidirectional traceability
verification, (iv) per-framework coverage analysis, and (v) gap
detection with categorical classification. The tool, deployment
scenarios, and generated results are publicly available.\footnote{%
\url{https://github.com/YOUR_USERNAME/ugaf-its-engine}}

We validated the engine against four architecturally distinct ITS
deployment archetypes selected to address the external validity
limitation identified in Section~\ref{subsec:threats}: (1)~a full
three-tier urban smart intersection mirroring the Dubai-inspired
scenario with 12~components across vehicle, edge/RSU, and cloud tiers;
(2)~a highway corridor ADS deployment with strong vehicle-tier autonomy
and minimal edge presence; (3)~a transit priority corridor operating
exclusively at the edge and cloud tiers with no vehicle AI; and
(4)~a rural intersection with a single-tier edge-only deployment
representing the minimum governance footprint. These four archetypes
span the deployment diversity described in Section~\ref{subsec:threats},
from maximum architectural complexity to minimum viable deployment.

Listing~\ref{lst:yaml_snippet} shows an abbreviated deployment
descriptor. The engine reads this descriptor, determines which tiers
are active, activates the applicable unified controls from the catalog
(Table~\ref{tab:uc_catalog}), computes the evidence backbone, and
verifies traceability completeness. A representative output fragment is
shown in Listing~\ref{lst:json_output}.

\begin{figure}[t]
\begin{lstlisting}[language={},basicstyle=\ttfamily\scriptsize,
  caption={Abbreviated YAML deployment descriptor for the transit
  priority corridor scenario.},label={lst:yaml_snippet},
  frame=single,xleftmargin=0.5em,numbers=none]
name: "Transit Priority Corridor"
components:
  - name: "Transit Signal Priority AI"
    tier: "T2"
    risk_level: "high"
    owner: "transit_authority"
  - name: "Schedule Optimization"
    tier: "T3"
    risk_level: "medium"
    owner: "transit_authority"
\end{lstlisting}
\end{figure}

\begin{figure}[t]
\begin{lstlisting}[language={},basicstyle=\ttfamily\scriptsize,
  caption={Representative JSON output fragment showing control
  activation and evidence metrics for the transit priority
  corridor.},label={lst:json_output},
  frame=single,xleftmargin=0.5em,numbers=none]
{
  "control_activation": {
    "total_activated": 12,
    "total_possible": 12,
    "activation_ratio": 100.0
  },
  "evidence_metrics": {
    "unified_artifacts": 18,
    "siloed_baseline": 37,
    "reduction_percentage": 51.4,
    "avg_reuse_factor": 2.89
  },
  "traceability": {
    "traceability_score": 100.0,
    "status": "COMPLETE"
  }
}
\end{lstlisting}
\end{figure}


% ============================================================================
% 2. NEW TABLE: MULTI-SCENARIO COMPARATIVE RESULTS
% ============================================================================
% Replace the current single-scenario Table 9 with this multi-scenario table.

\begin{table*}[t]
\centering
\caption{Multi-Scenario Computational Validation Results. Four
architecturally distinct ITS deployments evaluated by the UGAF-ITS
governance engine. 154 source obligations (55~ISO/IEC 42001 +
37~EU AI Act + 62~NIST AI RMF) consolidated into 12 unified controls.}
\label{tab:multi_scenario_results}
\renewcommand{\arraystretch}{1.15}
\setlength{\tabcolsep}{4pt}
\begin{tabular}{l c c c c}
\hline
\textbf{Metric}
  & \textbf{Urban Intersection}
  & \textbf{Highway Corridor}
  & \textbf{Transit Priority}
  & \textbf{Rural Intersection} \\
\hline
Active tiers
  & T1, T2, T3 & T1, T2, T3 & T2, T3 & T2 \\
Components
  & 12 & 7 & 5 & 2 \\
Stakeholders
  & 3 & 3 & 2 & 1 \\
\hline
Activated UCs
  & 12/12 & 12/12 & 12/12 & 11/12 \\
Unified artifacts
  & 20 & 20 & 18 & 12 \\
Siloed baseline
  & 37 & 37 & 37 & 37 \\
Evidence reduction
  & 45.9\% & 45.9\% & 51.4\% & 67.6\% \\
Avg.\ reuse factor
  & 2.80 & 2.80 & 2.89 & 2.92 \\
3-framework artifacts
  & 16 (80\%) & 16 (80\%) & 16 (89\%) & 11 (92\%) \\
\hline
ISO/IEC 42001 coverage
  & 94.5\% & 94.5\% & 94.5\% & 92.7\% \\
EU AI Act coverage
  & 91.9\% & 91.9\% & 91.9\% & 89.2\% \\
NIST AI RMF coverage
  & 88.7\% & 88.7\% & 88.7\% & 82.3\% \\
\textbf{Average coverage}
  & \textbf{91.7\%} & \textbf{91.7\%} & \textbf{91.7\%} & \textbf{88.1\%} \\
\hline
Traceability completeness
  & 100\% & 100\% & 100\% & 100\% \\
Coverage gaps
  & 13 & 13 & 13 & 19 \\
\hline
\end{tabular}
\begin{tablenotes}\footnotesize
\item \textit{Baseline}: siloed compliance modeled as 15~ISO + 12~EU + 10~NIST = 37 documents.
\textit{Rural gaps}: 6~tier-dependent (UC-11 inactive), 3~organizational, 3~regulatory, 7~context-setting.
All values computed by the UGAF-ITS governance engine (publicly available).
\end{tablenotes}
\end{table*}


% ============================================================================
% 3. REVISED ABSTRACT SENTENCES
% ============================================================================
% Replace the current single-scenario validation sentence with this:

% FIND in abstract:
%   "A Dubai-inspired smart intersection instantiation reports 95%
%    coverage of ISO/IEC 42001 Annex A controls, 92% coverage of scoped
%    EU AI Act high-risk obligations, and 88% coverage of NIST AI RMF
%    subcategories, while reducing evidence artifacts by 46%."
%
% REPLACE WITH:
%   "Computational validation across four architecturally distinct ITS
%    deployments---from full three-tier metropolitan intersections to
%    single-tier rural installations---reports 91.7\% average framework
%    coverage for three-tier deployments (ISO~94.5\%, EU AI Act~91.9\%,
%    NIST~88.7\%), evidence artifact reduction of 45.9\%--67.6\%
%    depending on architectural complexity, and 100\% bidirectional
%    traceability across all topologies. The governance engine and
%    deployment scenarios are publicly available."


% ============================================================================
% 4. REVISED SECTION VII-D (THREATS TO VALIDITY)
% ============================================================================
% Update the "External validity" paragraph. Replace:

% FIND:
%   "External validity: The intersection scenario reflects only one ITS
%    archetype. Highway corridors, transit-priority networks, and
%    low-connectivity deployments may shift which tiers dominate control
%    enforcement and evidence production."
%
% REPLACE WITH:
%   "External validity: We address the single-archetype limitation
%    through computational validation across four deployment archetypes
%    spanning three-tier urban intersections, highway corridors,
%    transit-priority networks, and single-tier rural intersections
%    (Table~\ref{tab:multi_scenario_results}). Results confirm that
%    evidence reduction remains stable (45.9\%--67.6\%) and traceability
%    is preserved at 100\% across architecturally diverse topologies.
%    Coverage varies from 88.1\% to 91.7\% depending on tier
%    composition, which is an honest architectural property: deployments
%    missing entire tiers legitimately require fewer controls.
%    Additional archetypes such as maritime corridors or airport surface
%    operations remain untested."


% ============================================================================
% 5. REVISED SECTION VII-C (VALIDATION RESULTS) — KEY PARAGRAPHS
% ============================================================================
% Add this paragraph after presenting Table multi_scenario_results:

% "The multi-scenario validation reveals four structural properties of
%  UGAF-ITS. First, governance requirements are architecturally
%  contingent: a rural single-tier deployment activates 11 of 12
%  controls and requires 12 artifacts, while a full three-tier
%  intersection activates all 12 controls and requires 20 artifacts.
%  The deactivated control (UC-11, Supplier AI Assessment) is
%  cloud-specific by design, confirming that the tier allocation model
%  correctly suppresses controls irrelevant to the deployment topology.
%  Second, evidence reduction remains stable across architectures:
%  45.9\% for three-tier deployments, rising to 67.6\% for single-tier
%  deployments where fewer tier-specific artifacts are needed while the
%  siloed compliance burden remains constant. Third, bidirectional
%  traceability is 100\% across all four scenarios, enforced by
%  construction through the crosswalk methodology. Fourth, coverage
%  gaps are systematically classifiable: 13 gaps in three-tier
%  deployments fall into organizational procedures (3), regulatory
%  workflows (3), and context-setting activities (7), while the rural
%  deployment adds 6 tier-dependent gaps where obligations target tiers
%  absent from the architecture."


% ============================================================================
% 6. SECTION TO ADD TO CONCLUSION
% ============================================================================
% Add before the final "Future work" sentence:

% "Computational validation across four ITS archetypes confirms that
%  the crosswalk methodology and evidence backbone produce stable
%  governance properties across architecturally diverse deployments.
%  The open-source governance engine enables reproducible evaluation
%  of arbitrary ITS architectures against the unified control catalog."


% ============================================================================
% END OF REVISION CONTENT
% ============================================================================
